% Three frames, one rotation, and how much of it gravity can give you.
\begin{tikzpicture}[font=\sffamily\scriptsize, x=1cm, y=1cm]

\tikzset{
  axline/.style={-Latex, line width=0.8pt},
  axlbl/.style={font=\sffamily\tiny, inner sep=1pt},
}

% ================================================================ the joint frame
\node[chlabel, anchor=south] at (-4.4,2.9) {the joint's frame};
\node[linkbody, minimum width=26mm, minimum height=6mm] at (-4.4,0.9) {};
\node[jointpin] (jp) at (-5.6,0.9) {};
\draw[axline, draw=linecol]      (-5.6,0.9) -- (-3.4,0.9);
\node[axlbl, anchor=north] at (-4.5,0.42) {x, along the link};
\draw[axline, draw=blue!55!black](-5.6,0.9) -- (-5.6,2.4) node[axlbl, anchor=south west]{z, the axis of rotation};
\draw[jointarc] (-4.7,1.5) arc (0:70:0.9);
\node[axlbl, anchor=west] at (-4.3,2.05) {the angle the encoder measures};

% ================================================================ the sensor frame
\node[chlabel, anchor=south] at (0.9,2.9) {the sensor's frame};
\node[draw=linecol, fill=hwcol, minimum width=14mm, minimum height=10mm,
      rotate=22, font=\sffamily\tiny] at (0.9,0.9) {shield};
\draw[axline, draw=linecol]      (0.9,0.9) -- ++(22:1.6) node[axlbl, right]{x'};
\draw[axline, draw=blue!55!black](0.9,0.9) -- ++(112:1.4) node[axlbl, above]{z'};
\node[axlbl, anchor=north, text width=32mm, align=center] at (0.9,0.1)
  {mounted however it fits: its axes are not the joint's};

% ================================================================ the rotation
\draw[bus, line width=1pt] (-3.0,0.9) -- node[lbl, above]{one rotation}
                                        node[lbl, below]{three angles} (-0.3,0.9);

% ================================================================ what gravity gives
\node[chlabel, anchor=south west] at (3.6,2.9) {what three static orientations give};
\foreach \i/\lab/\ang in {0/{flat}/270, 1/{on edge}/200, 2/{on end}/340}{
  \node[draw=linecol, fill=hwcol, minimum width=9mm, minimum height=7mm,
        font=\sffamily\tiny] (o\i) at (4.2+\i*1.7,1.7) {\lab};
  \draw[axline, draw=red!65!black] (4.2+\i*1.7,1.25) -- ++(\ang:0.85);
}
\node[axlbl, anchor=north, text width=54mm, align=center] at (5.9,0.55)
  {gravity, measured in the sensor's own axes, in three known attitudes};

% ================================================================ the caveat
\node[umlnote, text width=52mm, anchor=north west] at (-6.4,-0.9)
  {\textbf{Gravity fixes two angles well and the third one not at all.} Rotating
   the shield about the gravity vector does not change what gravity looks like,
   so that angle carries no information in a static measurement. It comes from
   the joint's own motion instead: turn the joint, and the sensor axis that
   reports the rate is the one aligned with the axis of rotation.};

\node[umlnote, text width=52mm, anchor=north west] at (0.4,-0.9)
  {The fitted rotation goes into the board description from chapter 4, with its
   evidence marked \emph{measured} rather than \emph{documented}, and with its
   residual stated. A rotation fitted to within half a degree is fine for
   vibration and contact detection and is not fine for anything that integrates,
   which is one more reason this volume does not integrate it.};

\node[chlabel, anchor=north west, text width=108mm] at (-6.4,-4.9)
  {The third frame, the link's, is the same as the joint's frame once the joint
   angle is known, which is exactly what makes an inertial unit on one link
   unable to give a joint angle on its own: it reports motion in space, and
   separating the joint's motion from the whole arm's needs either a second unit
   on the other side or the encoder that chapter 5 already provides.};
\end{tikzpicture}
